\section{Introduction}

\begin{displayquote}
I tried to break it to him gently [...] the only way to learn an unknown language is to interact with a native speaker [...] asking questions, holding a conversation, that sort of thing [...] If you want to learn the aliens' language, someone [...] will have to talk with an alien. Recordings alone aren't sufficient.

Ted Chiang, \emph{Story of Your Life}
\end{displayquote}

One of the main aims of AI is to develop agents that can cooperate with others to achieve goals \citep{wooldridge2009introduction}. Such coordination requires communication. If the coordination partners are to include humans, the most obvious channel of communication is natural language. Thus, handling natural-language-based communication is a key step toward the development of AI that can thrive in a world populated by other agents.

Given the success of deep learning models in related domains such as image captioning or machine
translation \citep[e.g.,][]{Sutskever:etal:2014,Xu:etal:2015}, it
would seem reasonable to cast the problem of training conversational
agents as an instance of supervised learning \citep{Vinyals:Le:2015}.
However, training on ``canned''
conversations does not allow learners to experience the interactive
aspects of communication. Supervised approaches, which focus on
the structure of language, are an excellent way to learn
general statistical associations between sequences of symbols. However, they do
not capture the functional aspects of communication, i.e., 
that humans use words to coordinate with others and 
make things happen \citep{Austin:1962,clark1996using,Wittgenstein:1953}.

This paper introduces the first steps of a research program based on
\textit{multi-agent coordination communication games}. These games
place agents in simple environments where they need to develop a language to 
coordinate and earn payoffs. Importantly, the agents start as blank slates,
but, by playing a game together, they can develop and bootstrap
knowledge on top of each others, leading to the emergence of a
language.

The central problem of our program, then, is the following: How do we
design environments that foster the development of a language that is
portable to new situations and to new communication partners (in
particular humans)?

We start  from the most basic challenge of using a language
in order to \emph{refer} to things in the context of a two-agent game.
We focus on two questions. First, whether \emph{tabula rasa} agents
succeed in communication. Second, what features of the environment 
lead to the development of codes resembling human
language. 

We assess this latter question in two ways. 
First, we consider whether the agents associate general
conceptual properties, such as broad object categories (as opposed to 
low-level visual properties), to the symbols they learn to use. Second, we 
examine whether the agents' ``word usage'' is partially interpretable by 
humans in an online experiment.

Other researchers have proposed communication-based environments for
the development of coordination-capable AI. 
Work in multi-agent systems has focused on the design of
pre-programmed communication systems to solve specific tasks (e.g.,
robot soccer, \citealt{stone1998towards}).  Most related to our work, \cite{sukhbaatar2016learning} and \cite{Foerster:etal:2016}
show that neural networks can evolve communication
in the context of games without a pre-coded protocol.
We pursue the same question, but further ask how we
can change our environment to make the emergent language more interpretable.

Others (e.g., the SHRLDU program of \citealt{Winograd:1971} or the game in \citealt{Wang:etal:2016}) propose building a communicating AI by putting humans in the loop from the very beginning. This approach has benefits but faces serious scalability issues, as active human intervention is required at each step. An attractive component of our game-based paradigm is that humans may be added as players, but do not need to be there all the time.

A third branch of research focuses on ``Wizard-of-Oz'' environments,
where agents learn to play games by interacting with a complex
scripted environment~\citep{Mikolov:etal:2015a}. This approach gives
the designer tight control over the learning curriculum, but imposes a
heavy engineering burden on developers. We also stress the
importance of the environment (game setup), but we focus on simpler environments with multiple agents that force them to get smarter by bootstrapping on top of each
other. 

We leverage ideas from work in linguistics, cognitive science and game
theory on the emergence of language 
\citep{Wagner:etal:2003,Skyrms:2010,
  crawford1982strategic,crawford:1998survey}. 
Our game is a variation
of Lewis' signaling game \citep{Lewis:1969}. There is a rich
tradition of linguistic and cognitive studies using similar setups
\citep[e.g.,][]{Briscoe:2002,Cangelosi:Parisi:2002,Spike:etal:2016,Steels:Loetzsch:2012}. What
distinguishes us from this literature is our aim to, eventually,
develop practical AI. This motivates our focus on more realistic input
data (a large collection of noisy natural images) and on trying to
align the agents' language with human intuitions.

Lewis' classic games have been studied extensively in game theory under the name of ``cheap talk''. These games have been used as models to study the evolution of
language both theoretically and experimentally
~\citep{crawford:1998survey,blume:1998experimental,crawford1982strategic}. A
major question in game theory is whether equilibrium actually occurs in a game as
convergence in learning is not guaranteed
\citep{fudenberg2014recency,roth1995learning}. And, if an equilibrium is reached, which one it will be
(since they are typically not unique). This is particularly true for
cheap talk games, which exhibit Nash equilibria in which precise language emerges, others where vague language
emerges and others where no language emerges at all
\citep{crawford1982strategic}. In addition, because in these games
language has no ex-ante meaning and only emerges in the context of the
equilibrium, some of the emergent languages may not be very natural.  Our results speak to both the convergence question and the
question of what features of the game cause the appearance of
different types of languages. Thus, our results are also of interest
to game theorists.

An evolutionary perspective has recently been advocated as a way
to mitigate the data hunger of traditional supervised approaches
\citep{Goodfellow:etal:2014, Silver:etal:2016}. This research confirms that 
learning can be
bootstrapped from \emph{competition} between agents. We focus,
however, on \emph{cooperation} between agents as a way to foster
learning while reducing the need for annotated data.
